% !PRINTED 511
\weberpage{511}

From this it follows that $\varrho$ must satisfy a quadratic equation, and hence that
\weq{16}{\varrho=\pm i\quad\text{or}\quad \varrho=\dfrac{\pm 1\pm\sqrt{-3}}{2}}
must hold, because these are the only non-real twelfth roots of unity which satisfy a quadratic equation, and $\varrho=\pm1$ would lead to $x=x'$, $y=y'$. Accordingly $\sqrt D$ must belong to a field which is determined by one of these irrationalities $\varrho$, that is, we must have
\weq{17}{D=-4m^2\quad\text{or}\quad D=-3m^2}
\gesperrt{In the first case} it follows from (15)
\weq{18}{\begin{aligned}
y'&=\pm 2mx,\\
2mx'&=\mp y,
\end{aligned}}
and hence, by (8):
\[
n=m^2(x^2+x'^2),
\]
and since $n$ was assumed to be odd, $m$ must also be odd; by (18) $y,y'$ are even, and because of (8) so also are $x,x'$, and from (12) it follows, since $\gamma$ and $\beta$ are even, that:
% [sic] German source claims $x,x'$ are both even, conflicting with $n=m^2(x^2+x'^2)$ for odd $n,m$
\[
M^4=(\alpha+\beta\omega)^6=\left(\dfrac{y+ixm}{2}\right)^6,
\]
and from (18):
\[
M'^4=(\alpha'+\beta'\omega)^6=\left(\dfrac{y'+ix'm}{2}\right)^6=-\left(\dfrac{y+ixm}{2}\right)^6,
\]
\weq{19}{M'^4=-M^4.}

Thus $M$ and $M'$ have not equal but opposite 4th powers.

In the second of the exceptional cases (17) there follows in the same way from (15):
\weq{20}{\begin{aligned}
\pm 2y'&=y-3mx,\\
\pm 2mx'&=y+mx.
\end{aligned}}

Thereupon it follows, if one takes the upper signs,
\weq{21}{m(x'y-y'x)=2n,}
and from this one concludes that $m$ must be odd; for if $m$ were even, then $y$ and $y'$ would have to be even, and by (21) $2n$ would be divisible by 4, and hence $n$ by 2. This case is therefore not to be considered further, if we assume an even $D$.

With this the following theorem is proved:

% the German breaks mid-word at "un-|geraden"; the English carries "odd divisor" over to 512
1. \gesperrt{Let $D\equiv 0\pmod 4$ be a negative discriminant, and let $x,y$ be two numbers without a common}
%CONT>>
